\section{Configuración MQTT y Seguridad}

\subsection{Parámetros del Broker MQTT}
En la Tabla~\ref{tab:config_broker} se detallan los parámetros de conexión del broker de mensajería implementado.

\begin{table}[H]
\centering
\small
\begin{tabularx}{\textwidth}{l p{4.5cm} X}
\toprule
\textbf{Parámetro} & \textbf{Valor} & \textbf{Descripción / Uso} \\
\midrule
\textbf{Software Broker} & Eclipse Mosquitto v2.0.x & Servidor de mensajería MQTT en GNU/Linux. \\
\textbf{Dirección IP} & \texttt{10.42.0.1} (WLAN) / \texttt{localhost} & IP del servidor para el ESP32 y para Node-RED. \\
\textbf{Puerto} & \texttt{1883/TCP} & Puerto estándar para tráfico MQTT en texto plano. \\
\textbf{Acceso Anónimo} & Desactivado (\texttt{allow\_anonymous false}) & Requiere autenticación por usuario y contraseña. \\
\textbf{Usuarios del Sistema} & \texttt{esp32\_user} y \texttt{dashboard} & Roles diferenciados para nodo de campo y servidor web. \\
\bottomrule
\end{tabularx}
\caption{Parámetros de configuración del broker MQTT.}
\label{tab:config_broker}
\end{table}

\subsection{Autenticación y Credenciales}
En el archivo \texttt{/etc/mosquitto/conf.d/local.conf} se establecieron las directivas de escucha y autenticación:
\begin{lstlisting}[style=custombash]
listener 1883 0.0.0.0
allow_anonymous false
password_file /etc/mosquitto/passwd
acl_file /etc/mosquitto/acl
\end{lstlisting}
Los usuarios se dieron de alta mediante \texttt{sudo mosquitto\_passwd /etc/mosquitto/passwd <usuario>}.

\subsection{Control de Acceso por Roles (ACL) y Firewall}
\begin{itemize}
    \item \proyectoBadge \textbf{Listas de Control de Acceso (ACL):} En \texttt{/etc/mosquitto/acl} se delimitaron los privilegios mínimos para aislar funciones:
    \begin{itemize}[noitemsep]
        \item \texttt{esp32\_user}: Permiso de escritura en \texttt{sensor/sht41/\#}, \texttt{actuador/led/estado} y \texttt{esp32/status}; permiso de lectura únicamente en \texttt{actuador/led/set}.
        \item \texttt{dashboard}: Permiso de lectura en \texttt{sensor/\#}, \texttt{actuador/led/estado} y \texttt{esp32/status}; permiso de escritura únicamente en \texttt{actuador/led/set}.
    \end{itemize}
    \item \proyectoBadge \textbf{Regla de Firewall UFW:} Se limitó el acceso al puerto 1883 exclusivamente a la subred del proyecto:
    \begin{lstlisting}[style=custombash]
sudo ufw allow from 10.42.0.0/24 to any port 1883 proto tcp
    \end{lstlisting}
\end{itemize}
